\chapter{Modalverben}
\label{cha:Modalverben}

Modalverben drücken Modalitäten wie Möglichkeit, Erlaubnis, Verpflichtung, Wunsch oder Vermutung aus.

\section{Übersicht der Modalverben}
\label{sec:ModalverbenUebersicht}

\begin{center}
\begin{tabular}{|l|c|c|c|c|}
\hline
\textbf{Verb} & \textbf{Bedeutung} & \textbf{Konjugation} & \textbf{Praeteritum} & \textbf{Perfekt} \\ \hline
dürfen & Erlaubnis, Möglichkeit & darf, darfst, darf, dürfen, dürft, dürfen & durfte & habe...dürfen \\ \hline
können & Fähigkeit, Möglichkeit & kann, kannst, kann, können, könnt, können & konnte & habe...können \\ \hline
müssen & Verpflichtung & muss, musst, muss, müssen, müsst, müssen & musste & habe...müssen \\ \hline
sollen & Empfehlung, Pflicht & soll, sollst, soll, sollen, sollt, sollten & sollte & habe...sollen \\ \hline
wollen & Wollen, Absicht & will, willst, will, wollen, wollt, wollen & wollte & habe...wollen \\ \hline
mögen & Sympathie, Wunsch & mag, magst, mag, mögen, mögt, mögen & mochte & habe...gemocht \\ \hline
\end{tabular}
\end{center}

\section{Grammatik aktiv: Modalverben im Satz}
\label{sec:ModalverbenSatz}

\begin{enumerate}
    \item Ich \underline{\hspace{1cm}} Deutsch lernen. (wollen) \quad \textit{(Lösung: will)}
    \item Du \underline{\hspace{1cm}} das nicht tun. (dürfen) \quad \textit{(Lösung: darfst)}
    \item Er \underline{\hspace{1cm}} gut schwimmen. (können) \quad \textit{(Lösung: kann)}
    \item Wir \underline{\hspace{1cm}} heute arbeiten. (müssen) \quad \textit{(Lösung: müssen)}
    \item Ihr \underline{\hspace{1cm}} pünktlich sein. (sollen) \quad \textit{(Lösung: sollt)}
    \item Sie \underline{\hspace{1cm}} Kaffee. (mögen) \quad \textit{(Lösung: mögen)}
\end{enumerate}

\section{Subjektive Bedeutung der Modalverben}
\label{sec:SubjektiveModalverben}

Modalverben können nicht nur objektive Modalitäten ausdrücken, sondern auch subjektive Einschätzungen:

\subsection{Weitergabe von Informationen (sollen/wollen)}
\begin{itemize}
    \item \textbf{sollen}: Man gibt Informationen wieder, die man gehört/gelesen hat. (Der Wahrheitsgehalt ist unsicher.)
        \begin{itemize}
            \item Er soll der beste Spieler sein. (= Man sagt, er sei der beste...)
            \item Sie soll in Paris leben. (= Ich habe gehört, dass sie in Paris lebt.)
        \end{itemize}
    \item \textbf{wollen}: Jemand behauptet etwas über sich selbst. (Distanzierung des Sprechers.)
        \begin{itemize}
            \item Er will der Beste sein. (= Er behauptet von sich, der Beste zu sein.)
            \item Sie will alles gewusst haben. (= Sie behauptet, alles gewusst zu haben.)
        \end{itemize}
\end{itemize}

\subsection{Ausdruck einer Vermutung}
\begin{itemize}
    \item \textbf{müssen/müsste}: Hohe Wahrscheinlichkeit (90-99\%)
        \begin{itemize}
            \item Er muss noch im Büro sein. (= Ich bin sicher, dass er noch im Büro ist.)
            \item Das muss stimmen. (= Es ist logisch, dass es stimmt.)
        \end{itemize}
    \item \textbf{dürfte}: Wahrscheinlichkeit (75\%)
        \begin{itemize}
            \item Er dürfte noch im Büro sein. (= Wahrscheinlich ist er noch im Büro.)
            \item Das dürfte falsch sein. (= Es ist wahrscheinlich falsch.)
        \end{itemize}
    \item \textbf{kann/könnte}: Möglichkeit (50\%)
        \begin{itemize}
            \item Das kann sein. (= Es ist möglich.)
            \item Das könnte stimmen. (= Es könnte vielleicht stimmen.)
        \end{itemize}
\end{itemize}

\section{Übungen: Subjektive Modalverben}
\label{sec:SubjektiveModalUebungen}

\begin{enumerate}
    \item Er \underline{\hspace{1cm}} der beste Spieler sein. (sollen/wollen - Gerücht) \quad \textit{(Lösung: soll)}
    \item Er \underline{\hspace{1cm}} der Beste sein. (wollen - Eigenbehauptung) \quad \textit{(Lösung: will)}
    \item Das \underline{\hspace{1cm}} stimmen. (müssen - logischer Schluss) \quad \textit{(Lösung: muss)}
    \item Er \underline{\hspace{1cm}} noch im Büro sein. (dürfte - wahrscheinlich) \quad \textit{(Lösung: dürfte)}
    \item Das \underline{\hspace{1cm}} falsch sein. (könnte - möglich) \quad \textit{(Lösung: könnte)}
    \item Sie \underline{\hspace{1cm}} krank sein. (sollen - Information) \quad \textit{(Lösung: soll)}
    \item Er \underline{\hspace{1cm}} das nicht gewusst haben. (wollen - Distanzierung) \quad \textit{(Lösung: will)}
\end{enumerate}

\section{Konjunktiv II Formen der Modalverben}
\label{sec:ModalKonjunktivII}

\begin{center}
\begin{tabular}{|l|c|c|}
\hline
\textbf{Infinitiv} & \textbf{Konjunktiv II} & \textbf{Bedeutung} \\ \hline
dürfen & dürfte & dürfte (höflich) \\ \hline
könnte & könnte & könnte (höflich) \\ \hline
müsste & müsste & müsste (höflich) \\ \hline
sollte & sollte & sollte (höflich) \\ \hline
wollte & wollte & wollte (unhöflich, nur „möchten" verwenden) \\ \hline
möchte & möchte & möchte (höflich) \\ \hline
\end{tabular}
\end{center}

\chapterend